% ================================================================== Chapitre 9
\chapter{Questions de cours et exercices corrigés}
\label{ch:exos}

\section*{Questions de cours (réponses courtes)}
\addcontentsline{toc}{section}{Questions de cours}

\begin{enumerate}[label=\textbf{Q\arabic*.},leftmargin=2.6em,itemsep=0.5em]
  \item \emph{Définir l'état standard de référence d'un élément. Quel est-il pour l'oxygène, le brome à $\SI{298}{K}$, le carbone ?}\\
  État standard du corps simple le plus stable à $T$ : \ce{O2(g)}, \ce{Br2(l)}, C graphite.
  \item \emph{Pourquoi $\Delta_fH^\circ(\ce{O2(g)}) = 0$ mais $\Delta_fH^\circ(\ce{O3(g)})\neq0$ ?}\\
  \ce{O2(g)} est l'ESR de l'oxygène : sa réaction de formation ne transforme rien. \ce{O3} n'est pas l'ESR.
  \item \emph{Énoncer la loi de Hess.}\\
  $\Delta_rH^\circ = \sum_i\nu_i\Delta_fH^\circ_i$ ; elle découle du fait que $H$ est une fonction d'état.
  \item \emph{Donner le signe de $\Delta_rS^\circ$ pour \ce{CaCO3(s) = CaO(s) + CO2(g)}.}\\
  $\Delta\nu_{\text{gaz}} = +1$ : $\Delta_rS^\circ>0$ (on trouve $+\SI{160}{J.K^{-1}.mol^{-1}}$).
  \item \emph{Établir la relation entre $\delta S_c$ et $\Delta_rG$ à $T$ et $P$ fixées.}\\
  $\dd H = \delta Q$, $\dd S = \delta Q/T + \delta S_c$, donc $\dd G = -T\delta S_c$ ; et $\dd G = \Delta_rG\,\dd\xi$. D'où $T\delta S_c = -\Delta_rG\,\dd\xi \geq 0$.
  \item \emph{Démontrer $\Delta_rG = \Delta_rG^\circ + RT\ln Q_r$.}\\
  Injecter $\mu_i = \mu_i^\circ + RT\ln a_i$ dans $\Delta_rG = \sum\nu_i\mu_i$.
  \item \emph{Énoncer la relation de Van~'t~Hoff et la démontrer dans l'approximation d'Ellingham.}\\
  $\dd\ln K^\circ/\dd T = \Delta_rH^\circ/RT^2$ ; dériver $\ln K^\circ = -\Delta_rH^\circ/RT + \Delta_rS^\circ/R$.
  \item \emph{Effet d'un ajout de gaz inerte à $P$ fixée sur \ce{N2 + 3H2 = 2NH3} ?}\\
  $Q_r\propto n_{\text{tot}}^{2}$ augmente : évolution dans le sens indirect (défavorable).
  \item \emph{Relier $\Delta_rG^\circ$ d'une réaction rédox aux potentiels standard.}\\
  $\Delta_rG^\circ = -\gamma F(E_1^\circ - E_2^\circ)$, d'où $\log K^\circ = \gamma(E_1^\circ - E_2^\circ)/0{,}06$ à $\SI{298}{K}$.
  \item \emph{Établir l'inégalité $\Delta G\leq W_{\text{élec}}$.}\\
  $\dd H = \delta Q + \delta W_{\text{élec}}$ (monobare), $\dd S = \delta Q/T + \delta S_c$, donc $\dd G = \delta W_{\text{élec}} - T\delta S_c$.
  \item \emph{Que valent $e$, $\Delta_rG$ et $Q_r$ pour une pile usée ?}\\
  $e = 0$, $\Delta_rG = 0$, $Q_r = K^\circ$ : la pile est à l'équilibre.
  \item \emph{Définir la capacité d'une pile. Convertir $\SI{1}{A.h}$ en coulombs.}\\
  $q_{\max} = \gamma F\xi_{\max}$ ; $\SI{1}{A.h} = \SI{3600}{C}$.
\end{enumerate}

\section*{Données pour les exercices}
\addcontentsline{toc}{section}{Données pour les exercices}
Les données thermodynamiques à $\SI{298}{K}$ sont rassemblées dans le tableau~\ref{tab:donnees} (chapitre~\ref{ch:formulaire}). On utilise l'approximation d'Ellingham, $R = \SI{8,314}{J.K^{-1}.mol^{-1}}$ et $F = \SI{96485}{C.mol^{-1}}$.

% ------------------------------------------------------------------
\exercice{Synthèse de l'ammoniac}{2}
\label{exo:nh3}
On étudie \ce{N2(g) + 3 H2(g) = 2 NH3(g)}.
\begin{enumerate}
  \item Prévoir le signe de $\Delta_rS^\circ$, puis calculer $\Delta_rH^\circ$ et $\Delta_rS^\circ$.
  \item Exprimer $\Delta_rG^\circ(T)$. Calculer la température d'inversion et $K^\circ(\SI{298}{K})$.
  \item Calculer $K^\circ$ à $\SI{723}{K}$ de deux façons.
  \item Dans un réacteur à $\SI{723}{K}$, les pressions partielles sont $P_{\ce{N2}} = \SI{50}{bar}$, $P_{\ce{H2}} = \SI{150}{bar}$, $P_{\ce{NH3}} = \SI{20}{bar}$. Dans quel sens le système évolue-t-il ?
  \item Justifier les conditions industrielles : $\SI{450}{\celsius}$, $\SI{200}{bar}$, catalyseur au fer, liquéfaction de \ce{NH3}.
\end{enumerate}

\correction
\begin{enumerate}
  \item $\Delta\nu_{\text{gaz}} = -2$ : $\Delta_rS^\circ<0$. $\Delta_rH^\circ = 2\Delta_fH^\circ(\ce{NH3}) = \SI{-91,8}{kJ.mol^{-1}}$ (corps simples dans leur ESR). $\Delta_rS^\circ = 2(192{,}8) - 191{,}6 - 3(130{,}7) = \SI{-198,1}{J.K^{-1}.mol^{-1}}$.
  \item $\Delta_rG^\circ(T) = -91{,}8\times10^3 + 198{,}1\,T$ (en $\mathrm{J\cdot mol^{-1}}$). $T_i = 91{,}8\times10^3/198{,}1 = \SI{463}{K}$. À $\SI{298}{K}$ : $\Delta_rG^\circ = \SI{-32,7}{kJ.mol^{-1}}$, $K^\circ = \exp(32{,}7\times10^3/(8{,}314\times298)) = \num{5,4e5}$.
  \item Directement : $\Delta_rG^\circ(723) = \SI{51,4}{kJ.mol^{-1}}$, $K^\circ = \num{1,9e-4}$. Par Van~'t~Hoff intégrée depuis $\SI{298}{K}$ : même valeur (voir chapitre~\ref{ch:vanthoff}).
  \item $Q_r = \dfrac{(P_{\ce{NH3}}/P^\circ)^2}{(P_{\ce{N2}}/P^\circ)(P_{\ce{H2}}/P^\circ)^3} = \dfrac{20^2}{50\times150^3} = \num{2,4e-6} < K^\circ$. Donc $\Delta_rG = RT\ln(Q_r/K^\circ) = \SI{-26}{kJ.mol^{-1}} < 0$ : évolution dans le \textbf{sens direct}, formation d'ammoniac.
  \item Réaction exothermique : $K^\circ$ diminue avec $T$, mais en dessous de $\SI{400}{\celsius}$ la réaction est trop lente même catalysée $\Rightarrow$ compromis vers $\SI{450}{\celsius}$. $\Delta\nu_{\text{gaz}}<0$ : une haute pression diminue $Q_r$ et favorise la synthèse. Le catalyseur accélère sans changer $K^\circ$. La liquéfaction retire \ce{NH3}, ce qui diminue $Q_r$ et permet de recycler \ce{N2} et \ce{H2}.
\end{enumerate}

% ------------------------------------------------------------------
\exercice{Effets thermiques : flamme et calorimétrie}{2}
\label{exo:gaz}
\textbf{A. Flamme dans le dioxygène pur.} On reprend la combustion \ce{CH4(g) + 2 O2(g) -> CO2(g) + 2 H2O(g)} ($\Delta_rH^\circ = \SI{-802,5}{kJ.mol^{-1}}$), mais avec un mélange stœchiométrique dans \ce{O2} pur, initialement à $\SI{298}{K}$.
\begin{enumerate}
  \item Calculer la température de flamme adiabatique et la comparer à celle obtenue dans l'air ($\approx\SI{2680}{K}$).
  \item La température réelle d'une flamme oxy-méthane est de l'ordre de $\SI{3000}{K}$. Commenter.
\end{enumerate}
\textbf{B. Mesure d'une enthalpie de réaction.} Dans un calorimètre de capacité thermique $C_{\text{cal}} = \SI{120}{J.K^{-1}}$, on mélange $\SI{100}{mL}$ d'acide chlorhydrique à $\SI{1,00}{mol.L^{-1}}$ et $\SI{100}{mL}$ de soude à $\SI{1,00}{mol.L^{-1}}$, tous deux à $\SI{20,0}{\celsius}$. On mesure une température finale de $\SI{25,8}{\celsius}$. On assimile les solutions à de l'eau ($c_{\text{eau}} = \SI{4,18}{J.g^{-1}.K^{-1}}$, $\rho = \SI{1}{g.mL^{-1}}$).
\begin{enumerate}[resume]
  \item Écrire la réaction et justifier qu'elle est totale.
  \item En déduire $\Delta_rH^\circ$. Comparer à la valeur tabulée calculée avec $\Delta_fH^\circ(\ce{H2O(l)}) = \SI{-285,8}{kJ.mol^{-1}}$, $\Delta_fH^\circ(\ce{HO-}) = \SI{-230,0}{kJ.mol^{-1}}$ et la convention $\Delta_fH^\circ(\ce{H+}) = 0$.
\end{enumerate}

\correction
\begin{enumerate}
  \item État final : \ce{CO2 + 2 H2O(g)}, $\sum n_jC_{p,m,j} = 37{,}1 + 2\times33{,}6 = \SI{104,3}{J.K^{-1}}$. $T_F = 298 + \dfrac{802{,}5\times10^3}{104{,}3} \approx \SI{8000}{K}$. Sans le diazote, il y a trois fois moins de matière à chauffer : la température calculée est trois fois plus élevée.
  \item La valeur calculée est absurde : à ces températures, \ce{CO2} et \ce{H2O} se dissocient (réactions très endothermiques qui « absorbent » l'énergie) et les $C_p$ augmentent. L'hypothèse « combustion totale, $C_p$ constants » n'est plus valable. Elle reste correcte qualitativement : la flamme dans \ce{O2} pur est bien plus chaude (d'où son usage en soudage et découpage).
  \item \ce{H+ + HO- -> H2O}, de constante $K^\circ = 1/K_e = 10^{14}$ : réaction quasi totale. $\xi_F = \SI{0,100}{mol}$.
  \item Calorimètre adiabatique, monobare : $\Delta H = 0 = \xi_F\Delta_rH^\circ + (m_{\text{eau}}c_{\text{eau}} + C_{\text{cal}})(T_F - T_I)$, avec $m_{\text{eau}}c_{\text{eau}} = 200\times4{,}18 = \SI{836}{J.K^{-1}}$ :
  \[ \Delta_rH^\circ = -\frac{(836 + 120)\times5{,}8}{0{,}100} = \SI{-55,4}{kJ.mol^{-1}}. \]
  Valeur tabulée : $\Delta_rH^\circ = -285{,}8 - (-230{,}0) - 0 = \SI{-55,8}{kJ.mol^{-1}}$. L'écart (moins de 1\,\%) vient des pertes thermiques et de l'assimilation des solutions à de l'eau pure.
\end{enumerate}

% ------------------------------------------------------------------
\exercice{Équilibre de Boudouard}{2}
\label{exo:boudouard}
On étudie \ce{C(graphite) + CO2(g) = 2 CO(g)}, réaction clé du haut fourneau.
\begin{enumerate}
  \item Calculer $\Delta_rH^\circ$ et $\Delta_rS^\circ$ ; commenter leurs signes. Calculer la température d'inversion.
  \item Calculer $K^\circ$ à $\SI{1000}{K}$ et à $\SI{800}{K}$.
  \item On place $\SI{1}{mol}$ de \ce{CO2} au contact d'un excès de graphite à $\SI{1000}{K}$, sous $P = \SI{1}{bar}$. Déterminer la composition à l'équilibre.
  \item Même question sous $P = \SI{10}{bar}$, puis sous $\SI{1}{bar}$ en présence de $\SI{4}{mol}$ de \ce{N2} (le gaz du haut fourneau provient de l'air). Interpréter.
\end{enumerate}

\correction
\begin{enumerate}
  \item $\Delta_rH^\circ = 2(-110{,}5) - (-393{,}5) = \SI{+172,5}{kJ.mol^{-1}}$ : endothermique. $\Delta_rS^\circ = 2(197{,}7) - 213{,}8 - 5{,}7 = \SI{+175,9}{J.K^{-1}.mol^{-1}}$ : positive, cohérente avec $\Delta\nu_{\text{gaz}} = +1$. $T_i = 172{,}5\times10^3/175{,}9 = \SI{981}{K}$.
  \item $K^\circ(1000) = \exp\left(-\dfrac{172{,}5\times10^3 - 1000\times175{,}9}{8{,}314\times1000}\right) = 1{,}50$ ; $K^\circ(800) = \num{8,4e-3}$. Réaction endothermique : $K^\circ$ croît avec $T$ (Van~'t~Hoff).
  \item Le graphite est en excès (activité 1). Avec $n_{\ce{CO2}} = 1-\xi$, $n_{\ce{CO}} = 2\xi$, $n_{\text{tot,gaz}} = 1+\xi$ :
  \[ K^\circ = \frac{(P_{\ce{CO}}/P^\circ)^2}{P_{\ce{CO2}}/P^\circ} = \frac{4\xi^2}{(1-\xi)(1+\xi)}\,\frac{P}{P^\circ} \;\Longrightarrow\; \xi_{\text{éq}} = \sqrt{\frac{K^\circ}{K^\circ + 4P/P^\circ}} = \sqrt{\frac{1{,}50}{5{,}50}} = \SI{0,52}{mol}. \]
  Composition : $\SI{0,48}{mol}$ de \ce{CO2}, $\SI{1,05}{mol}$ de \ce{CO}, soit $x_{\ce{CO}} = 0{,}69$.
  \item Sous $\SI{10}{bar}$ : $\xi_{\text{éq}} = \sqrt{1{,}50/41{,}5} = \SI{0,19}{mol}$ ($x_{\ce{CO}} = 0{,}32$) : $\Delta\nu_{\text{gaz}} > 0$, une hausse de pression défavorise la formation de CO. Avec $\SI{4}{mol}$ de \ce{N2} à $\SI{1}{bar}$ : $\dfrac{4\xi^2}{(1-\xi)(5+\xi)} = 1{,}50$ donne (résolution numérique) $\xi_{\text{éq}} = \SI{0,74}{mol}$ : la dilution par un inerte à $P$ fixée équivaut à une baisse de pression et favorise ici le sens direct. Le haut fourneau produit donc efficacement le réducteur CO à haute température et sous pression modérée.
\end{enumerate}

% ------------------------------------------------------------------
\exercice{Pile à combustible à hydrogène}{2}
\label{exo:pac}
Dans une pile à combustible à électrolyte acide, le dihydrogène est oxydé à une électrode et le dioxygène réduit à l'autre, avec formation d'eau liquide.
\begin{enumerate}
  \item Écrire les demi-équations, préciser l'anode et la cathode, et écrire la réaction de fonctionnement pour une mole de \ce{H2}.
  \item Calculer la tension à vide standard $e^\circ$ à $\SI{298}{K}$.
  \item Quelle énergie électrique maximale peut-on obtenir avec $\SI{1}{kg}$ de \ce{H2} ? Quelle capacité (en A$\cdot$h) cela représente-t-il ?
  \item On définit le rendement maximal $\eta = \Delta_rG^\circ/\Delta_rH^\circ$. Le calculer et l'interpréter.
  \item Calculer $\dd e^\circ/\dd T$, puis $e^\circ$ à $\SI{80}{\celsius}$ (température de fonctionnement usuelle).
  \item La cathode est alimentée par de l'air ($P_{\ce{O2}} = \SI{0,2}{bar}$) et l'anode par \ce{H2} sous $\SI{1}{bar}$. Calculer $e$ à $\SI{298}{K}$.
\end{enumerate}

\correction
\begin{enumerate}
  \item Anode ($\ominus$, oxydation) : \ce{H2 -> 2H+ + 2e^-}. Cathode ($\oplus$, réduction) : \ce{O2 + 4H+ + 4e^- -> 2H2O}. Bilan pour une mole de \ce{H2} : \ce{H2(g) + 1/2 O2(g) -> H2O(l)}, $\gamma = 2$.
  \item $\Delta_rG^\circ = \Delta_fG^\circ(\ce{H2O(l)}) = \SI{-237,1}{kJ.mol^{-1}}$, donc $e^\circ = \dfrac{237{,}1\times10^3}{2\times96485} = \SI{1,23}{V}$ (on retrouve $E^\circ(\ce{O2/H2O}) - E^\circ(\ce{H+/H2})$).
  \item $n = 1000/2{,}016 = \SI{496}{mol}$ ; $W_{\max} = n\times237{,}1 = \SI{118}{MJ} \approx \SI{33}{kWh}$. Capacité : $q_{\max} = 2Fn = \SI{9,6e7}{C} \approx \SI{2,7e4}{A.h}$.
  \item $\eta = 237{,}1/285{,}8 = 0{,}83$. Au mieux, 83\,\% de l'énergie de la réaction devient électrique ; le reste ($-T\Delta_rS^\circ$, avec $\Delta_rS^\circ<0$) est forcément cédé sous forme de chaleur. Ce rendement n'est pas limité par Carnot, contrairement à celui d'un moteur thermique.
  \item $e^\circ = -\dfrac{\Delta_rH^\circ - T\Delta_rS^\circ}{2F}$ donc $\dfrac{\dd e^\circ}{\dd T} = \dfrac{\Delta_rS^\circ}{2F}$, avec $\Delta_rS^\circ = 70{,}0 - 130{,}7 - \tfrac12(205{,}2) = \SI{-163,3}{J.K^{-1}.mol^{-1}}$ : $\dd e^\circ/\dd T = \SI{-0,85}{mV.K^{-1}}$. À $\SI{353}{K}$ : $e^\circ = 1{,}229 - 55\times0{,}85\times10^{-3} = \SI{1,18}{V}$.
  \item $Q_r = \dfrac{1}{(P_{\ce{H2}}/P^\circ)(P_{\ce{O2}}/P^\circ)^{1/2}}$, donc $e = e^\circ - \dfrac{0{,}059}{2}\log Q_r = 1{,}229 + \dfrac{0{,}059}{4}\log 0{,}2 = \SI{1,22}{V}$. L'effet de la dilution du dioxygène est faible.
\end{enumerate}

% ------------------------------------------------------------------
\exercice{Degrés d'oxydation du cuivre}{2}
\label{exo:cuivre}
Données : $E^\circ(\ce{Cu^{2+}/Cu}) = \SI{0,34}{V}$, $E^\circ(\ce{Cu^{2+}/Cu+}) = \SI{0,16}{V}$.
\begin{enumerate}
  \item Calculer $E^\circ(\ce{Cu+/Cu})$.
  \item Écrire la réaction de dismutation de \ce{Cu+} et calculer $\Delta_rG^\circ$ et $K^\circ$.
  \item On prépare une solution de \ce{Cu+} à $\SI{0,010}{mol.L^{-1}}$. Que se passe-t-il ? Calculer la concentration résiduelle en \ce{Cu+}.
\end{enumerate}

\correction
\begin{enumerate}
  \item \ce{Cu^{2+} + 2e^- = Cu} est la somme de \ce{Cu^{2+} + e^- = Cu+} et \ce{Cu+ + e^- = Cu} : $-2F(0{,}34) = -F(0{,}16) - F\,E^\circ(\ce{Cu+/Cu})$, donc $E^\circ(\ce{Cu+/Cu}) = 2\times0{,}34 - 0{,}16 = \SI{0,52}{V}$.
  \item \ce{2 Cu+ = Cu^{2+} + Cu(s)} ($\gamma = 1$) : \ce{Cu+} est oxydant du couple \ce{Cu+/Cu} ($\SI{0,52}{V}$) et réducteur du couple \ce{Cu^{2+}/Cu+} ($\SI{0,16}{V}$). $\Delta_rG^\circ = -F(0{,}52 - 0{,}16) = \SI{-34,7}{kJ.mol^{-1}}$, $K^\circ = 10^{0{,}36/0{,}06} = 10^6$.
  \item $K^\circ$ est grand : dismutation quasi totale, formation de cuivre métallique et de \ce{Cu^{2+}} à environ $\SI{0,005}{mol.L^{-1}}$. Alors $\dfrac{[\ce{Cu^{2+}}]c^\circ}{[\ce{Cu+}]^2} = 10^6$ donne $[\ce{Cu+}] = \sqrt{0{,}005\times10^{-6}} = \SI{7e-5}{mol.L^{-1}}$. \ce{Cu+} n'est pas stable en solution aqueuse : la règle du gamma le montre aussi, car l'oxydant \ce{Cu+} est au-dessus du réducteur \ce{Cu+} sur l'échelle des potentiels.
\end{enumerate}

% ------------------------------------------------------------------
\exercice{Pile bouton argent--zinc}{2}
\label{exo:bouton}
Une pile bouton de montre fonctionne selon \ce{Zn(s) + Ag2O(s) -> ZnO(s) + 2 Ag(s)} dans un électrolyte de potasse concentrée.

\noindent Données : $\Delta_fG^\circ(\ce{ZnO(s)}) = \SI{-320,5}{kJ.mol^{-1}}$ ; $\Delta_fG^\circ(\ce{Ag2O(s)}) = \SI{-11,2}{kJ.mol^{-1}}$ ; $M(\ce{Ag2O}) = \SI{231,7}{g.mol^{-1}}$ ; $M(\ce{Zn}) = \SI{65,4}{g.mol^{-1}}$.
\begin{enumerate}
  \item Écrire les demi-équations en milieu basique ; identifier l'anode et la cathode.
  \item Calculer la tension à vide standard. Pourquoi la tension à vide reste-t-elle constante pendant toute la décharge ?
  \item La pile contient $\SI{0,50}{g}$ d'oxyde d'argent (réactif limitant). Calculer sa capacité en mA$\cdot$h et la masse minimale de zinc nécessaire.
  \item Une montre consomme $\SI{2}{\micro A}$ en moyenne. Estimer la durée de vie de la pile.
\end{enumerate}

\correction
\begin{enumerate}
  \item Anode ($\ominus$) : \ce{Zn + 2HO- -> ZnO + H2O + 2e^-}. Cathode ($\oplus$) : \ce{Ag2O + H2O + 2e^- -> 2Ag + 2HO-}. Bilan avec $\gamma = 2$.
  \item $\Delta_rG^\circ = -320{,}5 - (-11{,}2) = \SI{-309,3}{kJ.mol^{-1}}$, $e^\circ = \dfrac{309{,}3\times10^3}{2\times96485} = \SI{1,60}{V}$. Tous les constituants du bilan sont des solides purs ($a = 1$) : $Q_r = 1$ en permanence, donc $e = e^\circ$ \emph{tout au long} de la décharge (les ions \ce{HO-} et l'eau s'éliminent du bilan). C'est une qualité recherchée pour l'horlogerie.
  \item $n(\ce{Ag2O}) = 0{,}50/231{,}7 = \SI{2,16}{mmol} = \xi_{\max}$ ; $q_{\max} = 2F\xi_{\max} = \SI{416}{C} = \SI{116}{mA.h}$. Zinc : $\SI{2,16}{mmol}\times65{,}4 = \SI{141}{mg}$ au minimum. Énergie : $q_{\max}e^\circ = \SI{667}{J}$.
  \item $\Delta t = q_{\max}/i = 416/(2\times10^{-6}) = \SI{2,1e8}{s} \approx 6{,}6$ ans (en pratique un peu moins, à cause de l'autodécharge).
\end{enumerate}
